% ─────────────────────────────────────────────────────────────────────────────
% Math typesetting — the packages and operators every booklet needs.
%
% Wired into both the EN and FA templates (booklets/templates/*/main.tex).
% amsthm is loaded for its \theoremstyle machinery only; the actual
% definition/theorem/example boxes are the tcolorbox environments in
% configs/env.tex, not amsthm's \newtheorem counters.
% ─────────────────────────────────────────────────────────────────────────────
\usepackage{amsmath}
\usepackage{amssymb}
\usepackage{mathtools}
\usepackage{amsthm}
\usepackage{bm}

\DeclarePairedDelimiter{\floor}{\lfloor}{\rfloor}
\DeclarePairedDelimiter{\ceil}{\lceil}{\rceil}
\DeclarePairedDelimiter{\abs}{\lvert}{\rvert}
\DeclarePairedDelimiter{\norm}{\lVert}{\rVert}
\DeclareMathOperator*{\argmin}{arg\,min}
\DeclareMathOperator*{\argmax}{arg\,max}

% Number sets
\newcommand{\Naturals}{\ensuremath{\mathbb{N}}}
\newcommand{\Integers}{\ensuremath{\mathbb{Z}}}
\newcommand{\Rationals}{\ensuremath{\mathbb{Q}}}
\newcommand{\Reals}{\ensuremath{\mathbb{R}}}
\newcommand{\Complex}{\ensuremath{\mathbb{C}}}

% Probability
\newcommand{\Prob}{\ensuremath{\mathbb{P}}}
\newcommand{\Expect}{\ensuremath{\mathbb{E}}}
\DeclareMathOperator{\Var}{Var}

% Set notation
\newcommand{\set}[1]{\{#1\}}
\newcommand{\setbuilder}[2]{\{#1 \mid #2\}}
\newcommand{\symdiff}{\ensuremath{\triangle}}

% ── Unicode math symbols typed directly into prose ───────────────────────────
% Authors keep typing ℝ, π, ε, → etc. straight into chapter text instead of
% $\mathbb{R}$, $\pi$, $\varepsilon$, $\to$ — harmless under xelatex (native
% Unicode), a hard compile error under pdflatex. Map the ones that have
% actually shown up so the en/ pipeline doesn't choke on them; \ensuremath
% makes each mapping work whether it lands in text or math mode.
\ifdefined\XeTeXversion
\else
  \usepackage{newunicodechar}
  \newunicodechar{ℝ}{\ensuremath{\mathbb{R}}}
  \newunicodechar{ℤ}{\ensuremath{\mathbb{Z}}}
  \newunicodechar{ℕ}{\ensuremath{\mathbb{N}}}
  \newunicodechar{ℚ}{\ensuremath{\mathbb{Q}}}
  \newunicodechar{ℂ}{\ensuremath{\mathbb{C}}}
  \newunicodechar{π}{\ensuremath{\pi}}
  \newunicodechar{ε}{\ensuremath{\varepsilon}}
  \newunicodechar{δ}{\ensuremath{\delta}}
  \newunicodechar{ζ}{\ensuremath{\zeta}}
  \newunicodechar{Θ}{\ensuremath{\Theta}}
  \newunicodechar{→}{\ensuremath{\to}}
  \newunicodechar{²}{\textsuperscript{2}}
  \newunicodechar{³}{\textsuperscript{3}}
  \newunicodechar{¹}{\textsuperscript{1}}
  \newunicodechar{ⁿ}{\textsuperscript{\textit{n}}}
  \newunicodechar{ᵖ}{\textsuperscript{\textit{p}}}
  \newunicodechar{₀}{\textsubscript{0}}
  \newunicodechar{₁}{\textsubscript{1}}
\fi
